\begin{table*}[!tbp]
\centering\small
\caption{Regole di costruzione degli otto vicinati sociali. Il parametro $k$ dei grafi R4 e R8 controlla la densità di campionamento.}
\label{tab:app-graphs}
\begin{tabularx}{\textwidth}{@{}lX@{}}
\toprule
Vicinato & Collegamenti tra particelle \\
\midrule
Global & Ogni particella ha accesso al migliore dell'intero sciame. \\
Ring & Ogni nodo è collegato ai due adiacenti nell'anello. \\
Von Neumann & Quattro adiacenze ortogonali su una griglia toroidale. Con $N_p=64$ la griglia è $8\times 8$. \\
Stella & Un hub scelto casualmente è collegato a tutti gli altri nodi. \\
Random-4 (R4) & Archi Bernoulli con probabilità $4/(N_p-1)$. Eventuali componenti vengono unite mediante ponti. \\
Random-8 (R8) & Archi Bernoulli con probabilità $8/(N_p-1)$ e stessa regola di connessione. \\
Small-world & Anello di grado quattro con probabilità di ricollegamento $0.15$ per arco. \\
Scale-free & Attaccamento preferenziale con due archi per nuovo nodo, a partire da tre nodi interamente connessi. \\
\bottomrule
\end{tabularx}
\end{table*}
